\chapter{Text and fonts}
\label{ch:text}

Lettering is most of what a laser is used for: names, part numbers, signs, serial
numbers, batch codes. MeerK40t handles text in two fundamentally different ways,
and choosing the right one decides how the job looks.

\section{Two kinds of text}

\begin{description}[leftmargin=1.4em,style=nextline]
  \item[Bitmap text] An ordinary text element (\btn{Text} in the \panel{Tools}
        bar, or \cmd{text WORKSHOP} in the console). It behaves like a picture of
        text: fine for engraving, and it can also be converted into paths.
  \item[Vector text] Text drawn as paths from a font file, created with the
        \emph{linetext} tool. Every stroke is a real vector, so it can be cut,
        scaled to any size without losing crispness, welded, and given a stroke
        width for a bolder look.
\end{description}

\noindent Which to use:

\begin{itemize}
  \item \textbf{Cutting letters out} (a stencil, a sign, letters with holes):
        vector text, assigned to a \mkseries{Cut} operation.
  \item \textbf{Engraving crisp small lettering}: vector text in an
        \mkseries{Engrave} operation. Thin, clean, repeatable.
  \item \textbf{Large filled letters} on a plaque: either an \mkseries{Engrave}
        operation with a \emph{hatch effect} (Chapter~\ref{ch:ops}) or a
        \mkseries{Raster} operation, which fills the letters by scanning them.
  \item \textbf{Text that will be edited in a picture editor}, or where the
        exact typography matters more than the geometry: bitmap text, burned as
        a raster.
\end{itemize}

\section{Choosing a font}

MeerK40t can use three families of font:

\begin{center}\small
\begin{tabular}{@{}L{2.4cm}L{4.0cm}L{7.8cm}@{}}
\toprule
\textbf{Type} & \textbf{Extension} & \textbf{Character} \\
\midrule
Hershey fonts & \file{.jhf} & The classic single-stroke engraving fonts. Thin,
even, ideal for small lettering and for cutting. \\
Autocad shape fonts & \file{.shx} & Single-stroke engineering lettering; crisp
and utilitarian. \\
TrueType fonts & \file{.ttf} & Ordinary desktop fonts with outlines. Great for
display lettering, but thin fonts vanish at small sizes when cut. \\
\bottomrule
\end{tabular}
\end{center}

\noindent Open the \btn{Font-Manager} (from \menu{Edit>Settings>Font-Manager}, or
the preferences window's \btn{Fonts} button). It shows an information panel, a
\emph{Font-Work-Directory} with a picker for your own fonts, the list of
\emph{System font locations} MeerK40t searches, the list of fonts it has found
with a preview, and the buttons \btn{Import}, \btn{Delete}, \btn{Refresh}. The
\emph{Goto a font-source} combo jumps to the well-known collections --- the
Hershey font sets and the Autocad SHX fonts. Importing reports how many fonts
succeeded and how many failed, which is worth reading: one bad file in a batch
does not stop the rest.

\begin{tipbox}[Pick fonts that survive the laser]
A font with hairlines will burn into a thin scorch mark, and a font with tight
counters (the enclosed parts of \cmd{e}, \cmd{a}, \cmd{o}) will fill in when it
is cut small. Test your intended font at the size you intend to use --- a line of
text on scrap costs seconds and settles the argument completely.
\end{tipbox}

\section{Working with text elements}

Select a text element and choose \btn{Text properties} (or \btn{LineText
properties} for vector text). The property panel offers:

\begin{center}\small
\begin{tabular}{@{}L{3.4cm}L{10.8cm}@{}}
\toprule
\textbf{Control} & \textbf{What it does} \\
\midrule
\emph{Content} & The text itself. Multi-line content is supported. \\
\emph{Left / Center / Right} & Alignment for multi-line text. \\
font size buttons & Larger and smaller by one step. Hold \key{Shift} or
\key{Ctrl} for bigger steps; right-click resets to the default. \\
character-gap buttons & Increase or decrease the kerning between characters. \\
line-distance buttons & Increase or reduce the space between lines. \\
\btn{Weld} & Merges overlapping characters into one outline --- essential for
script fonts and for anything where two letters touch, since otherwise they cut
as two pieces that fall apart. \\
font list & Double-click a font to apply it there and then; the preview below
updates as you browse. \\
\btn{Insert symbol} & Right-click the text box for \btn{Paste} and
\btn{Insert symbol}, which opens a glyph picker listing each character with its
Unicode value, ASCII code and preview --- the practical way to add
\textdegree{}, \textmu{}, \texteuro{} or a non-Latin character. \\
\bottomrule
\end{tabular}
\end{center}

\noindent If a font in an imported design is not installed, the panel says
\emph{Font missing} and names it. Install the font, or switch the text to a font
you have --- do not begin a burn with a missing-font warning on screen, because
what you are looking at is not what will be cut.

\section{Text that changes between burns: wordlists}

A wordlist turns a text element into a template. Write a placeholder in the text,
such as \cmd{{NAME}}, and MeerK40t substitutes the current wordlist entry at
burn time. The index then advances by one, so the next burn of the same job
carries the next value --- which is what makes a single project produce a hundred
numbered parts.

\subsection*{The Wordlist Editor}

Open \menu{Edit>Settings>Wordlist-Editor}, or the \panel{Wordlist} pane with its
\btn{Edit} button. Each entry has a \emph{Name}, a type and an \emph{Index}. The
editor's controls include:

\begin{itemize}
  \item \emph{Filter} to find an entry by name;
  \item \emph{Start Index for CSV-based data} and \emph{Start Index for field} to
        control where counting begins;
  \item \btn{Jump} to position the index directly;
  \item \btn{Add Text} and the list tools (add, edit, delete entries and their
        values), and a tab containing the how-to text for the feature.
\end{itemize}

\noindent Besides the lists you create, several entries exist automatically and
are useful in labels and serial plates: \cmd{version}, \cmd{date}, \cmd{time},
\cmd{op_device}, \cmd{op_speed}, \cmd{op_power}, \cmd{op_dpi} and
\cmd{op_passes}. A date-stamped test piece costs nothing to produce when the text
element contains \cmd{{date}} and \cmd{{time}}.

\subsection*{Counting without a computer in the loop}

The \panel{Wordlist} pane's \btn{Next} and \btn{Prev} buttons step the whole
list forwards or backwards --- left-click for pages and right-click for single
entries. Between jobs, one click in that pane is often all the bookkeeping a
production run needs. The element context menu has matching
\btn{Increase Wordlist-Reference} and \btn{Decrease Wordlist-Reference} entries
for adjusting a single reference, and \cmd{{NAME#+1}} in the text reaches
forward within the list while writing.

\subsection*{CSV data}

A list can come from a file: the header row of a CSV becomes the entry names, and
the rows become the values, with the \emph{Start Index for CSV-based data} setting
deciding where to begin. That is the mechanism for engraving names onto a stack of
gifts, or part numbers onto a tray of brackets, from a spreadsheet.

\begin{walkthrough}{Number ten tags without touching the design}
\begin{enumerate}[label=\arabic*.,leftmargin=1.4em]
  \item Design one tag with a text element reading \cmd{No. {SERIAL}}.
  \item Open the Wordlist Editor, add a wordlist called \cmd{SERIAL} and enter
        ten values: 001 to 010.
  \item Set the index to the first value.
  \item Make a cutting sheet of ten tags, all using that same text element.
  \item Burn the job. Every tag carries the next value in sequence, and the index
        is left ready for the next sheet.
  \item Check the index in the \panel{Wordlist} pane before re-running anything;
        if the job is repeated, the numbering continues rather than restarting.
\end{enumerate}
\end{walkthrough}

\section{Turning text into geometry}

Three conversions appear on the text element's context menu, and they exist to
solve three different problems:

\begin{itemize}
  \item \btn{Convert to path} --- the text becomes paths you can edit node by
        node. Do this when you must adjust a single letter.
  \item \btn{Convert to vector text} --- a bitmap text element is re-drawn as
        vector text using the font of the same name. Do this when the design
        arrived as outlines and you want stroke-based lettering.
  \item \btn{Trace bitmap} --- vectorises the rendered text. Rarely the best
        option, but it works on text of unknown origin
        (Chapter~\ref{ch:images}).
\end{itemize}

\begin{tryit}{Make a stencil}
\begin{enumerate}[label=\arabic*.,leftmargin=1.4em]
  \item Create vector text with your chosen word or initials, sized so the
        lettering is at least 20\unit{mm} tall.
  \item Choose a stencil-friendly font, or use \btn{Weld} on a script font so
        the letters hold together.
  \item Convert it to a path and use \btn{Break Subpaths} to separate the
        letters.
  \item Add small connector bridges with \btn{Tab Edit} or by drawing short lines
        across the cut, so the counters (the middles of \cmd{A}, \cmd{O}, \cmd{R})
        do not fall out.
  \item Assign everything to a \mkseries{Cut} operation, position it, and
        cut it in card or thin plywood.
  \item Adjust the bridges and re-cut. The second attempt is always the one you
        keep --- which is the whole point of testing in card first.
\end{enumerate}
\end{tryit}
